%! Author = murfe
%! Date = 27.08.2026
\chapter{Einleitung}
\label{ch:einleitung}
\section{Motivation}
\label{sec:motivation}
Ein Zahlungsauftrag, der nach einem Verbindungsabbruch versehentlich mehrfach ausgeführt wird, belastet einen Kunden doppelt.
Ein Compliance-Report, dessen Erstellung durch den Absturz des ausführenden Prozesses unbemerkt verloren geht, fehlt zum
gesetzlich vorgeschriebenen Stichtag.
So unterschiedlich beide Fälle wirken, beruhen sie auf derselben grundlegenden Anforderung an ein Background-Job-Processing-System:
Ein Job darf im Fehlerfall weder mehrfach noch gar nicht wirken und vor allem nicht verloren gehen.
Was sich zunächst als Frage der internen Konsistenz eines Systems liest, kann damit unmittelbar finanzielle und regulatorische
Konsequenzen haben, die weit über das zugrundeliegende technische Problem hinausreichen.\\
In modernen Softwaresystemen stellt die Verarbeitung von Hintergrundjobs einen zentralen Bestandteil vieler Backend-Architekturen
dar \parencite[Kap.~10]{kleppmann_designing_2017}.
Zahlreiche Geschäftsprozesse, darunter die Generierung von Reports, der Import großer Datenmengen, periodische Abrechnungsläufe
oder das Versenden asynchroner Benachrichtigungen, erfordern eine Verarbeitung, die entkoppelt vom synchronen Request-Response-Zyklus stattfindet.
Insbesondere in industrienahen und sicherheitskritischen Anwendungen müssen solche Prozesse nicht nur funktional korrekt ablaufen,
sondern darüber hinaus zuverlässig, nachvollziehbar und fehlertolerant sein.
Aber auch in nicht-kritischen Systemen ist Zuverlässigkeit zum einen von Bedeutung für die Nutzerfreundlichkeit der Anwendung,
zum anderen aber auch um Produktivitäts- und Gewinneinbußen in Produktivsystemen in Folge von Fehlern zu vermeiden
\parencite[Kap.~1 Abschnitt \glqq How Important Is Reliability?\grqq]{kleppmann_designing_2017}.

\section{Problemstellung}
\label{sec:problemstellung}
In der Praxis können bei der Verarbeitung von Hintergrundjobs jederzeit Teilausfälle auftreten.
Prozesse können während der Verarbeitung unerwartet terminieren, Datenbankverbindungen können kurzzeitig unterbrochen werden, identische
Anfragen infolge von clientseitigen Wiederholungsversuchen mehrfach an das System übermittelt werden oder
Ergebnisse bzw. Antworten können verloren gehen \parencite[Kap.~8]{kleppmann_designing_2017}.
Insbesondere Netzwerkpartitionierungsfehler tauchen dabei in der Praxis als häufige Ausfallursachen in Cloud-Systemen auf
\parencite{alquraan_analysis_2018}, wodurch die im Folgenden beschriebenen Konsistenzprobleme keine seltenen Randfälle,
sondern einen im laufenden Betrieb regelmäßig zu erwartenden Fall darstellen.
Ohne geeignete Mechanismen führen derartige Szenarien zu inkonsistenten Systemzuständen oder dem Verlorengehen einer Nachricht.
Eine naheliegende Reaktion auf letzteren Fall besteht darin, betroffene Operationen automatisch zu wiederholen.
Werden jedoch viele zuvor fehlgeschlagene oder unterbrochene Operationen ohne geeignete Steuerung gleichzeitig erneut versucht,
etwa weil eine Datenbank nach einem kurzzeitigen Ausfall wieder verfügbar wird und sämtliche zuvor blockierten Anfragen
nahezu zeitgleich erneut auf sie zugreifen, kann dies selbst zu einer erheblichen, kurzfristigen Lastspitze führen, die
das gerade erst wiederhergestellte System erneut überlastet.
Dieses in der Praxis als \textit{Thundering-Herd-Problem} bezeichnete Phänomen zeigt, dass naive Fehlerbehandlungsstrategien
ein bereits bestehendes Problem unter Umständen sogar verschärfen können \parencite[Kap.~4]{newman_building_2026}.

Bestehende Frameworks und Bibliotheken, etwa Quartz Scheduler \parencite{terracotta_inc_hrsg_quartz_2026}
oder Spring Batch \parencite{broadcom_inc_hrsg_spring_2026}, stellen umfangreiche Funktionalitäten für die Hintergrundverarbeitung bereit.
Hierzu zählen unter anderem Scheduling, Wiederholungsmechanismen sowie Verfahren zur Fehlerbehandlung und Wiederaufnahme.
Sie verfolgen dabei jedoch jeweils eigene Architekturansätze und abstrahieren die zugrunde liegenden Mechanismen weitgehend
hinter ihren Programmierschnittstellen.
Dadurch eignen sie sich nur eingeschränkt, um systematisch zu untersuchen, welche Architekturentscheidungen für bestimmte
Konsistenz- und Verarbeitungszusicherungen tatsächlich erforderlich sind.

Die beschriebene Problematik ist von hoher praktischer und industrieller Relevanz.
Nahezu jedes Backend-System, das über einfache \gls{crud}-Operationen hinausgeht, muss Aufgaben zuverlässig im Hintergrund
verarbeiten und dabei konsistente Datenbestände gewährleisten \parencite[Kap.~4]{richardson_microservices_2018}.
Fehler in diesem Bereich führen zu Datenverlust, inkonsistenten Geschäftszuständen sowie unnötigem manuellem Eingriff und
können insbesondere in regulierten Domänen, wie beispielsweise der Luftfahrt, Compliance-Verstöße nach sich ziehen.
Auch im Kontext von Cloud-nativen Architekturen ist dieses Thema relevant.
Horizontale Skalierung, Container-Neustarts und kurzzeitige Netzwerkunterbrechungen gehören dort zum regulären Betriebsverhalten
und müssen bereits beim Entwurf eines Systems berücksichtigt werden.
Daher bilden die in dieser Arbeit behandelten Konzepte auch die Grundlage für den zuverlässigen Betrieb von Backend-Systemen in Cloud-Umgebungen.\\
Zur Umsetzung einer solchen zuverlässigen Hintergrundverarbeitung hat sich in der Praxis der Einsatz dedizierter Message
Broker wie RabbitMQ oder Azure Service Bus etabliert, die eine robuste Entkopplung zwischen Auftragserstellung und -verarbeitung
ermöglichen \parencite[Kap.~11]{kleppmann_designing_2017}.
Dieser Ansatz erfordert jedoch den zusätzlichen Betrieb und die Überwachung einer eigenständigen Infrastrukturkomponente
neben der ohnehin für die fachlichen Daten benötigten Datenbank.
Zusätzlich braucht es eine Konsistenzsicherung, etwa mittels des Transactional-Outbox-Patterns, um Datenbank und Broker
konsistent zu halten \parencite[Kap.~8 Abschnitt \glqq Pattern: Transactional Outbox\grqq]{newman_building_2026}.
Für viele klein- und mittelgroße Systeme steht dieser zusätzliche Betriebsaufwand in keinem angemessenen Verhältnis zum
tatsächlichen Bedarf.
Daraus ergibt sich als leitende Überlegung dieser Arbeit, ob sich die angestrebten Zuverlässigkeitsgarantien auch allein
mit einer ohnehin bereits vorhandenen Datenbank erreichen lassen, ohne den Betrieb einer zusätzlichen Broker-Infrastruktur.

Obwohl die zugrundeliegenden Konzepte wie Fehlertoleranz, Recovery, Idempotenz, Nebenläufigkeitskontrolle und Transaktionsmanagement
in der Literatur umfassend beschrieben werden, fehlen Arbeiten, die diese Aspekte zu einer durchgängigen datenbankbasierten
Architektur für Background-Job-Processing zusammenführen und deren Verhalten unter definierten Fehlerszenarien systematisch evaluieren.
Ebenso implementieren bestehende Frameworks diese Mechanismen häufig als interne Abstraktionen, wodurch deren Zusammenspiel und Auswirkungen
auf die Zuverlässigkeit des Gesamtsystems für Entwickler nur eingeschränkt nachvollziehbar sind.
Daraus ergibt sich die dieser Arbeit zugrundeliegende Hauptforschungsfrage:\\
Welche Architekturmechanismen und Entwurfsentscheidungen ermöglichen es, in einem Java/Spring-basierten Job-Processing-System
definierte Konsistenzzusicherungen unter typischen Teilausfallszenarien nachweisbar einzuhalten?\\
Zur Beantwortung dieser Frage werden folgende Teilfragen untersucht:
\begin{enumerate}[label=TF\arabic*:, ref=TF\arabic*]
    \item \label{item:tf1} Welche Zusicherungen muss das System gewährleisten, damit Jobs stets
    korrekt und ohne Inkonsistenzen verarbeitet werden?
    \item \label{item:tf2} Welchen Beitrag leisten Idempotenz, Retry-Strategien und Recovery-Mechanismen
    zur Robustheit unter definierten Fehlerszenarien?
    \item \label{item:tf3} Wie müssen Transaktionen und Locking-Strategien genutzt werden, damit
    bei paralleler Arbeit mehrerer Worker keine Doppelverarbeitung oder inkonsistente Zustände entstehen?
    \item \label{item:tf4} Wie lässt sich sicherstellen, dass Jobs nach einem unerwarteten Prozessabbruch korrekt
    \item wiederaufgenommen werden, ohne dass Aufträge verloren gehen?
\end{enumerate}

\section{Zielsetzung}
\label{sec:zielsetzung}
Ziel dieser Arbeit ist der Entwurf, die prototypische Implementierung und die Evaluation
eines fehlertoleranten Job-Processing-Backends auf Basis einer relationalen Datenbank in Kombination mit Java und Spring.
Dabei soll herausgearbeitet werden, welche Architekturmechanismen erforderlich sind, damit Hintergrundjobs auch bei Teilausfällen
zuverlässig und korrekt verarbeitet werden und definierte Konsistenzzusicherungen eingehalten werden.

Konkret verfolgt die Arbeit folgende Teilziele:
\begin{enumerate}[label=TZ\arabic*:, ref=TZ\arabic*]
    \item \label{item:tz1} Definition von Zusicherungen: Es sollen präzise Korrektheitsbedingungen
    formuliert werden, die das System einhalten muss, unabhängig davon, welche Fehler auftreten.
    Dazu zählt insbesondere, dass kein Job doppelt erfolgreich abgeschlossen wird und dass kein Job dauerhaft verloren geht.
    \item \label{item:tz2} Architekturentwurf: Es soll eine Architektur entworfen werden, die die hergeleiteten Zusicherungen umsetzt.
    Architekturentscheidungen sollen explizit dokumentiert und begründet werden.
    \item \label{item:tz3} Prototypische Implementierung: Die entworfene Architektur soll als
    funktionsfähiger Prototyp umgesetzt werden, der die definierten Zusicherungen
    unter realistischen Bedingungen einhält.
    \item \label{item:tz4} Evaluation unter Fehlerbedingungen: Es soll nachgewiesen werden, dass das
    System seine Zusicherungen auch bei gezielt herbeigeführten Teilausfällen nicht
    verletzt.
\end{enumerate}

\section{Hypothesen}
\label{sec:hypothesen}
Auf Basis der Forschungsfrage und der Zielsetzung werden folgende Hypothesen formuliert,
deren Überprüfung im Rahmen der Evaluation erfolgt.
\begin{enumerate}[label=H\arabic*:, ref=H\arabic*]
    \item \label{item:h1} Durch die Kombination eines persistenten Statusmodells mit transaktionsbasiertem
    Locking und Idempotenz-Mechanismen lässt sich ein System entwerfen, das auch bei
    unerwarteten Prozessabbrüchen und Teilausfällen keine Jobs verliert und keine doppelten Ergebnisse
    erzeugt.
    \item \label{item:h2} Retry-Strategien mit exponentiellem Backoff in Verbindung mit einer
    Fehlerklassifikation (transient vs. permanent) erhöhen die Robustheit des Systems
    gegenüber kurzzeitigen Infrastrukturausfällen, ohne dabei die definierten Zusicherungen
    zu verletzen.
    \item \label{item:h3} Die definierten Zusicherungen lassen sich durch automatisierte Tests mit gezielter
    Fehlerinjektion reproduzierbar überprüfen und nachweisen.
\end{enumerate}

\section{Aufbau der Arbeit}
\label{sec:aufbau}
\Cref{ch:grundlagen} legt die theoretischen Grundlagen der Arbeit.
Es führt zunächst die Grundkonzepte asynchroner Verarbeitung und des Background-Job-Processings ein, bevor die für
Nebenläufigkeit und Fehlertoleranz zentralen Konzepte wie \gls{acid}, Synchronisationsverfahren, Fehlerklassifikation,
Retry-Strategien, Recovery und Idempotenz vorgestellt werden.
Das Kapitel schließt mit einer Einordnung des Stands der Technik, aus der sich die in \cref{sec:problemstellung}
skizzierte Lücke ergibt, die die vorliegende Arbeit füllt.\\
\Cref{ch:methodik} beschreibt das methodische Vorgehen der Arbeit.
Dazu zählt das Vorgehen bei der Literaturrecherche, die Erfassung der Anforderungen sowie die daraus schrittweise Ableitung
von Zusicherungen, als auch die Methodik der späteren Evaluation.
Dort werden insbesondere das Vorgehen zur Operationalisierung der Zusicherungen und die verwendeten Test- und Nachweismethoden betrachtet.\\
\Cref{ch:anforderungen} legt die Anforderungen an das zu entwickelnde System dar.
Nach einer Beschreibung des Systemkontexts wird das betrachtete Fehlermodell explizit abgegrenzt.
Darauf aufbauend werden die Zusicherungen Z-1 bis Z-9 hergeleitet, die als verbindlicher Maßstab für Entwurf und Evaluation dienen.\\
\Cref{ch:umsetzung} beschreibt den Entwurf und die prototypische Implementierung des Systems.
Es begründet die zentralen Architekturentscheidungen, darunter die Verwendung einer relationalen Datenbank als Job-Store, das
Zustandsmodell sowie die Mechanismen für Nebenläufigkeitskontrolle, Retry und Recovery und dokumentiert deren spezifische
Umsetzung in Java und Spring.\\
\Cref{ch:evaluation} evaluiert das implementierte System.
Ausgehend von der in \cref{tab:testmatrix} dokumentierten Testmatrix werden gezielt Fehler in das System injiziert und
die resultierenden Systemzustände gegen die definierten Prüfkriterien geprüft.
Die Ergebnisse werden anschließend zur Bewertung der in \cref{sec:hypothesen} formulierten Hypothesen herangezogen.\\
\Cref{ch:diskussion} diskutiert die gewonnenen Erkenntnisse sowie die Grenzen und Einschränkungen der Arbeit, insbesondere %TODO: Hier am Ende nochmal prüfen
hinsichtlich des in \cref{sec:abgrenzungFehlermodell} bewusst eingeschränkten Fehlermodells.
Des Weiteren gibt das Kapitel einen Ausblick auf mögliche Erweiterungen, besonders hinsichtlich Skalierbarkeit und Performance unter höherer Last.
